\section{Galilean spacetime and the geometry of reference frames}

Classical mechanics distinguishes the time of an event from its spatial
location.  The distinction can be stated without first choosing an observer,
a spatial origin, or coordinate axes.  We begin with the affine formulation
of Prastaro, \emph{Geometry of PDEs and Mechanics}, Definition 3.1 and
Remarks 3.1--3.3, pp.~307--309.  The affine assumptions made here are
stronger than merely having a time-fibered manifold: they also supply the
flat affine connection used below to identify inertial motion.

\subsection{Events, absolute time, and simultaneous space}

An affine space modeled on $V$ is a set on which $V$ acts freely and
transitively by translations: for any two points $p,q$ there is a unique
$v\in V$ with $p+v=q$.  Let $M$ be a four-dimensional affine space modeled
on $V$, and let $T$ be a one-dimensional affine space modeled on $V_T$.  An
\emph{event} is a point $p\in M$.  A point of an affine space is not itself
a vector: for $p,q\in M$, the difference $q-p\in V$ is defined, but a
position vector for $p$ requires a chosen origin.

Absolute time is a surjective affine map
\begin{equation}\label{eq:gal-time}
 \tau:M\longrightarrow T,\qquad p\longmapsto\tau(p).
\end{equation}
Its linear part $D\tau:V\to V_T$ is independent of $p$.  Thus, for two
events $p,q$,
\begin{equation}\label{eq:gal-time-difference}
 \tau(q)-\tau(p)=D\tau(q-p).
\end{equation}
The kernel
\begin{equation}\label{eq:gal-spatial-vectors}
 S=\ker D\tau\subset V,\qquad \dim S=3,
\end{equation}
is the vector space of spatial displacements.  The simultaneous space at
$t\in T$ is the fiber $M_t=\tau^{-1}(t)$.  If $p,q\in M_t$, then
\begin{equation}\label{eq:gal-simultaneity}
 D\tau(q-p)=\tau(q)-\tau(p)=0,
 \qquad q-p\in S.
\end{equation}
Conversely, if $p\in M_t$ and $\mathbf s\in S$, then
$\tau(p+\mathbf s)=t$.  Therefore $M_t$ is an affine space modeled on $S$.
The fibers at different times are distinct; the time map alone does not
identify a point of one fiber with a point of another as being at rest.

A positive-definite inner product $g:S\times S\to\mathbb R$ gives the
spatial distance between simultaneous events:
\begin{equation}\label{eq:gal-distance}
 d_t(p,q)^2=g(q-p,q-p),\qquad p,q\in M_t.
\end{equation}
This is a metric on each spatial fiber, not a positive-definite spacetime
metric.  In Prastaro's notation the affine model space is $\mathbf M$, the
spatial model space is $\mathbf S$, and the constant Euclidean form is
$\underline g$; our $V$, $S$, and $g$ denote those respective objects.
The vertical bundle is $\ker D\tau\subset TM$; for this affine model it
is naturally $M\times S$, with $g$ acting on its fibers.

Choose a time origin and time unit, identifying $T$ with $\mathbb R$.
Writing $\sigma=d\tau$, an adapted chart $(x^0,x^1,x^2,x^3)$ may be
chosen with $x^0=\tau$ up to a constant.  Then
\begin{equation}\label{eq:gal-adapted}
 \sigma\!\left(\frac{\partial}{\partial x^0}\right)=1,
 \qquad
 \sigma\!\left(\frac{\partial}{\partial x^i}\right)=0,
 \qquad i=1,2,3.
\end{equation}
The three coordinate vectors $\partial/\partial x^i$ are tangent to the
simultaneous fiber.  They need not be orthonormal.  For an orthonormal
Cartesian basis $\mathbf e_i$ of $S$, $g(\mathbf e_i,\mathbf e_j)=\delta_{ij}$,
and the spatial line element in Eq.~\eqref{eq:gal-distance} becomes
\begin{equation}\label{eq:gal-cartesian-metric}
 d\ell^2=\delta_{ij}\,dx^i dx^j.
\end{equation}
This is the spatial Euclidean expression used by Zipfel, \emph{Modeling and
Simulation of Aerospace Vehicle Dynamics}, 1st ed., \S2.1.1.  It is a
specialization of the metric on simultaneous space, rather than a metric
including time.

\subsection{Worldlines and velocities}

A section of a map $\pi:E\to B$ is a map $s:B\to E$ such that
$\pi\circ s=\operatorname{id}_B$; the definition applies in particular to
bundle projections.
\begin{definition}[Motion]
A motion is a smooth section $\gamma:T\to M$ of the time map:
\begin{equation}\label{eq:gal-motion}
 \tau\circ\gamma=\operatorname{id}_T.
\end{equation}
\end{definition}
For $t\in T$, the derivative is a linear map
\begin{equation}\label{eq:gal-derivative-map}
 D\gamma(t):T_tT\longrightarrow T_{\gamma(t)}M.
\end{equation}
After choosing the time origin and unit, identify $T_tT$ with $\mathbb R$
and $T_{\gamma(t)}M$ with the translation space $V$.  Evaluating the map
in Eq.~\eqref{eq:gal-derivative-map} on $1\in\mathbb R$ gives
\begin{equation}\label{eq:gal-tangent-lift}
 \dot\gamma(t)=D\gamma(t)(1)\in V,\qquad
 t\longmapsto(\gamma(t),\dot\gamma(t))\in TM\simeq M\times V.
\end{equation}
The first expression is the vector part; the second is the tangent lift,
which also records its base event.  Prastaro denotes the motion by $m$
and its lift by $\dot m$ (Remark 3.3, pp.~308--309).
The chain rule applied to Eq.~\eqref{eq:gal-motion} gives
\begin{equation}\label{eq:gal-velocity-normalization}
 D\tau\bigl(\dot\gamma(t)\bigr)
 =\frac{d}{dt}\bigl(\tau(\gamma(t))\bigr)=1.
\end{equation}
Thus the spacetime velocity lies in the affine hyperplane
\begin{equation}\label{eq:gal-velocity-bundle}
 I=\{v\in V:D\tau(v)=1\},\qquad
 J_D(M)=M\times I\subset TM.
\end{equation}
This is Prastaro's space of velocities of motions (Definition 3.3); its
fiber $I$ is affine, not the vector space $S$.  Indeed, if $v,w\in I$,
then $D\tau(v-w)=1-1=0$, so $v-w\in S$; conversely $v+s\in I$ for
$s\in S$.  For two motions at the same time,
\begin{equation}\label{eq:gal-relative-velocity}
 D\tau\bigl(\dot\gamma_1-\dot\gamma_2\bigr)=1-1=0,
 \qquad \dot\gamma_1-\dot\gamma_2\in S.
\end{equation}
Their relative velocity is consequently a spatial vector.  This makes precise
why a spatial velocity requires a reference motion even though a spacetime
worldline has a tangent independently of that choice.

\subsection{Reference frames and inertial motion}

A reference frame may be described by a congruence of worldlines: through
each event there passes one worldline representing a point at rest in that
frame.  Equivalently, locally, it is a time-preserving identification of
$M$ with $T\times\Sigma$, where $\Sigma$ labels the worldlines.  A frame
therefore identifies spatial points at different times; it is more than a
single worldline and does not by itself require numerical coordinates.
Prastaro's frame $\psi$ supplies such a congruence (Definition 3.4 and
Proposition 3.1, p.~309).  Zipfel's frame is a rigid set of mutually resting
points (\S2.1.4, pp.~20--21), so a general congruence need not be a Zipfel
frame: the latter additionally imposes rigidity.

More explicitly, Prastaro's $\psi:T\times M\to M$ sends an event $p$
to the event on its frame worldline at time $t$; the maps obey
$\tau(\psi(t,p))=t$, $\psi(\tau(p),p)=p$, and
$\psi(t',\psi(t,p))=\psi(t',p)$.  Thus $p\sim q$ if $q$ lies on the
same worldline as $p$, and the set $\Sigma=M/{\sim}$ labels frame
worldlines.  The splitting maps an event to
$(\tau(p),[p])\in T\times\Sigma$.  For a smooth frame its normalized
velocity field is
\begin{equation}\label{eq:gal-frame-field}
 U(p)=\left.\frac{d}{dt}\psi(t,p)\right|_{t=\tau(p)},
 \qquad D\tau(U(p))=1.
\end{equation}
The spatial velocity of a motion relative to this frame is
\begin{equation}\label{eq:gal-frame-relative}
 \mathbf v_{\mathrm{rel}}(t)=\dot\gamma(t)-U(\gamma(t))\in S,
\end{equation}
by Eq.~\eqref{eq:gal-relative-velocity}.  A frame need not be rigid or
inertial to define this subtraction.

The affine structure of $M$ supplies a flat connection.  A worldline is
inertial when it is an affine straight line, equivalently
\begin{equation}\label{eq:gal-inertial}
 \frac{d^2\gamma}{dt^2}=0,
\end{equation}
where the derivative is taken using that affine structure.  A standard global
inertial frame consists of parallel such worldlines.  Neither the time map
nor the spatial metric alone specifies this inertial structure on a general
fibered manifold.  An accelerated or rotating rigid frame can be used to
label events, but its coordinate expressions of the equations of motion
contain the familiar noninertial terms.

\subsection{The covariant derivative and ordinary mechanics}

Prastaro's absolute differential $\nabla X$ is a covector-valued vector
field: at each event it maps a direction $Y$ to $\nabla_YX$ (Definition 3.5
and Proposition 3.2, p.~310).  In bundle language, $X$ is a
section of $TM\to M$ and $DX$ takes values in $T(TM)$; the connection
extracts the vertical change, identified with a vector in the original
tangent fiber.  This makes derivatives of vectors at different points
meaningful without treating those tangent spaces as identical by assumption.

For coordinates $x^\alpha$, $\alpha=0,1,2,3$, define the connection
coefficients by
\begin{equation}\label{eq:gal-connection-basis}
 \nabla_{\partial/\partial x^\alpha}
 \frac{\partial}{\partial x^\gamma}
 =\Gamma^\beta{}_{\alpha\gamma}
 \frac{\partial}{\partial x^\beta}.
\end{equation}
For $X=X^\beta\partial/\partial x^\beta$ and
$Y=Y^\alpha\partial/\partial x^\alpha$, the product rule and
Eq.~\eqref{eq:gal-connection-basis} give
\begin{align}\label{eq:gal-covariant-coordinate}
 \nabla_YX
 &=Y^\alpha\left(
 \frac{\partial X^\beta}{\partial x^\alpha}
 \frac{\partial}{\partial x^\beta}
 +X^\gamma\Gamma^\beta{}_{\alpha\gamma}
 \frac{\partial}{\partial x^\beta}\right)\\
 &=Y^\alpha\left(
 \frac{\partial X^\beta}{\partial x^\alpha}
 +\Gamma^\beta{}_{\alpha\gamma}X^\gamma\right)
 \frac{\partial}{\partial x^\beta}.\nonumber
\end{align}
This is Prastaro's formula with his coordinate derivatives written as
explicit partial derivatives.  Along a motion, its acceleration is
$\nabla_{\dot\gamma}\dot\gamma$, defined as the covariant derivative along
the curve.  In components,
\begin{equation}\label{eq:gal-curve-acceleration}
 a^\beta=\frac{d^2x^\beta}{dt^2}
 +\Gamma^\beta{}_{\alpha\gamma}
 \frac{dx^\alpha}{dt}\frac{dx^\gamma}{dt}.
\end{equation}
For the affine connection in inertial Cartesian coordinates the connection
coefficients vanish.  With $x^0=t$, Eq.~\eqref{eq:gal-curve-acceleration}
then gives $a^0=0$ and $a^i=d^2x^i/dt^2$.

This is the setting of the familiar particle equation in Goldstein-style
classical mechanics.  Choose an inertial frame, an origin worldline, and
fixed Cartesian axes; write the particle's spatial position as
$\mathbf r=x^i\mathbf e_i$.  For constant mass $m$ and applied spatial force
$\mathbf F$, Newton's equation becomes
\begin{equation}\label{eq:gal-newton-goldstein}
 m\mathbf a=m\frac{d^2\mathbf r}{dt^2}=\mathbf F.
\end{equation}
The familiar $\mathbf r$ depends on the chosen origin worldline, whereas
the acceleration is the intrinsic acceleration for the chosen flat affine
connection.  In curvilinear or moving coordinates, the connection terms in
Eq.~\eqref{eq:gal-curve-acceleration} account for the changing coordinate
basis; they do not by themselves introduce a new physical force.

\subsection{Frames and coframes on a vector bundle}

Choquet-Bruhat distinguishes a tangent-space frame at one point from a
frame field on an open set (\S\S4.1--4.2, pp.~3--4).
The distinction extends to any rank-$n$ vector bundle $\pi:E\to B$.
A local frame over $U\subset B$ is a collection of smooth sections
$e_1,\ldots,e_n$ such that $e_i(x)$ is a basis of the fiber $E_x$ at
every $x\in U$.  The dual bundle $E^*$ has the coframe
$\theta^1,\ldots,\theta^n$, uniquely defined by
\begin{equation}\label{eq:gal-dual-frame}
 \theta^i_x(e_j(x))=\delta^i{}_j.
\end{equation}
For $v\in E_x$ and $\alpha\in E_x^*$, their components and reconstruction
are
\begin{equation}\label{eq:gal-frame-reconstruction}
 v^i=\theta^i_x(v),\quad v=v^i e_i(x),\qquad
 \alpha_i=\alpha(e_i(x)),\quad \alpha=\alpha_i\theta^i_x.
\end{equation}
A frame field selects a basis separately in each fiber, with smooth dependence
on the base point.  To compare vectors in distinct fibers geometrically,
however, a frame alone does not select a preferred transport: a connection
supplies that additional rule.  A metric permits us to say that frames are
orthonormal and to interpret changes between them as rotations.

For $E=TB$, the frame vectors are tangent vector fields and the coframe
consists of one-forms.  In coordinates $(x^1,\ldots,x^n)$, the natural
frame and coframe are
\begin{equation}\label{eq:gal-natural-frame}
 e_i=\frac{\partial}{\partial x^i},\qquad
 \theta^i=dx^i,\qquad
 dx^i\!\left(\frac{\partial}{\partial x^j}\right)=\delta^i{}_j.
\end{equation}
Not every frame field is a coordinate frame.  For example, on the plane
with the origin removed, the orthonormal polar frame and its dual are
\begin{equation}\label{eq:gal-polar-frame}
 e_r=\frac{\partial}{\partial r},\quad
 e_\phi=\frac{1}{r}\frac{\partial}{\partial\phi},\qquad
 \theta^r=dr,\quad\theta^\phi=r\,d\phi.
\end{equation}
Here $[e_r,e_\phi]=-r^{-1}e_\phi\ne0$, whereas coordinate vector fields
commute.  Thus a moving frame need not arise from a coordinate system.

A global frame is a local frame defined on all of $B$.  It exists precisely
when $E$ is isomorphic to the product bundle $B\times\mathbb R^n$: the
isomorphism sends $(x,(a^1,\ldots,a^n))$ to $a^i e_i(x)$.  For $E=TB$
this property is called parallelizability.  Choquet-Bruhat notes that not
every manifold has a global tangent frame, while every smooth orientable
three-manifold is parallelizable (p.~4); the latter is a theorem, not a
consequence of orientability in arbitrary dimensions.  For example, the
hairy-ball theorem rules out a global tangent frame on $S^2$.  This does not
contradict orientability of $S^2$, nor does the three-manifold theorem say
that every vector bundle over a three-manifold is trivial.

\subsection{Moving axes in engineering mechanics}

In the affine Galilean model, the spatial bundle is $M\times S$, so a fixed
orthonormal basis of $S$ gives a global spatial frame.  A frame attached to a
rotating vehicle may instead be chosen along its worldline.  Let
$\mathbf e_i(t)$ be a differentiable right-handed orthonormal basis of $S$,
and write a spatial vector along the motion as
$\mathbf v(t)=v^i(t)\mathbf e_i(t)$.  The affine identification of the
spatial fibers lets us differentiate these vectors in the common space $S$:
\begin{equation}\label{eq:gal-moving-product}
 \frac{d\mathbf v}{dt}
 =\frac{dv^i}{dt}\mathbf e_i
 +v^i\frac{d\mathbf e_i}{dt}.
\end{equation}
Differentiating $g(\mathbf e_i,\mathbf e_j)=\delta_{ij}$ shows that
$\Omega_{ij}=g(\mathbf e_i,d\mathbf e_j/dt)$ is skew-symmetric.  In
three-dimensional oriented Euclidean space there is a unique angular
velocity vector $\boldsymbol\omega$ such that
$d\mathbf e_i/dt=\boldsymbol\omega\times\mathbf e_i$.  Substitution in
Eq.~\eqref{eq:gal-moving-product} gives the transport formula
\begin{equation}\label{eq:gal-moving-transport}
 \frac{d\mathbf v}{dt}
 =\left(\frac{d\mathbf v}{dt}\right)_{\!\mathrm{axes}}
 +\boldsymbol\omega\times\mathbf v,\qquad
 \left(\frac{d\mathbf v}{dt}\right)_{\!\mathrm{axes}}
 =\frac{dv^i}{dt}\mathbf e_i.
\end{equation}
The first term differentiates components in the moving axes; the second
accounts for their rotation.  This is the geometric basis of the rotational
derivative used in vehicle dynamics.  A spatial basis field and a
reference-frame congruence are different choices: axes alone do not specify
which worldlines represent rest, and the congruence alone does not choose
axes.

\subsection{Axes, components, and Zipfel's notation}

At a given time let $\{\mathbf e_i\}$ and $\{\mathbf f_i\}$ be two
orthonormal bases of $S$.  Define the matrix $R$ by
\begin{equation}\label{eq:gal-basis-map}
 \mathbf f_j=\mathbf e_i R^i{}_{j},\qquad R^i{}_{j}=g(\mathbf e_i,\mathbf f_j).
\end{equation}
For the same spatial vector $\mathbf v$,
\begin{equation}\label{eq:gal-passive}
 \mathbf v=\mathbf e_i v_E^i=\mathbf f_j v_F^j
 =\mathbf e_i R^i{}_{j}v_F^j,
 \qquad [\mathbf v]_E=R[\mathbf v]_F.
\end{equation}
Here the basis changes and the geometric vector does not: this is a passive
component transformation.  Orthogonality gives
$[\mathbf v]_F=R^{\mathsf T}[\mathbf v]_E$.  By contrast, an active rotation
maps $\mathbf v$ to a generally different vector while a chosen basis is held
fixed.  The same numerical matrix can represent either operation only after
its domain, codomain, and basis conventions have been specified.

Zipfel's brackets denote components in the indicated coordinate system, not
a new physical vector (\S2.1.5, pp.~22--23).  His $\mathbf s_{AB}$ denotes
position of point $A$ with respect to $B$, and $[s_{AB}]^E$ is its column
of components in Earth coordinates.  For simultaneous event worldlines
$\gamma_A$ and $\gamma_B$, our translation is
\begin{equation}\label{eq:gal-zipfel-displacement}
 \mathbf s_{AB}(t)=\gamma_A(t)-\gamma_B(t)\in S,
 \qquad [s_{AB}]^E=[\mathbf s_{AB}]_E.
\end{equation}
Zipfel's $\mathbf v_B^I$ denotes velocity of point $B$ relative to inertial
frame $I$, whereas $[v_B^I]^E$ expresses that \emph{same} relative velocity
in Earth axes.  The reference frame determining velocity and the axes used
to display its components are independent choices.  His $\boldsymbol\omega^{BI}$
is angular velocity of frame $B$ relative to $I$; $[\omega^{BI}]^I$ is its
component column in inertial axes.  These uppercase labels identify frames
and axes, not tensor indices.  A transformation matrix must be defined by
its direction before it can be identified with $R$ or $R^{\mathsf T}$ in
Eq.~\eqref{eq:gal-passive}.
